\originalchapter{电荷与电场}\label{ch:charge}
\originalsection{带电现象与守恒}
塑料棒摩擦后吸引轻物体, 是因为电子在材料之间转移, 使一方有过量电子而另一方缺电子. 电荷是物体参与电磁作用的性质, 其数值记为$q$, 单位库仑(C). 电子电荷为$-e$, 质子为$+e$, $e=1.602176634\times10^{-19}\,\mathrm C$. 普通宏观物体几乎中性, 中性是正负电荷的代数和接近零, 不代表内部没有电荷.

对不与外界交换电荷的系统, 总电荷守恒. 金属有可移动的电子, 带电后可发生电荷重新分布; 绝缘体中的电荷在当前时间尺度内较难作长距离移动. “导体”和“绝缘体”是材料及时间尺度的描述, 不能把导体中的正负总电荷必然为零当成定义. 接地允许与地球交换电荷, 所以被接地物体自身的电荷可以变化, 但物体加地球的系统仍守恒.

\originalsection{库仑力与叠加}
位置矢量$\rr=x\uv x+y\uv y+z\uv z$从选定原点指向物体, $\uv x,\uv y,\uv z$为三条坐标轴正向的单位箭头. 两矢量相加逐分量相加; 点积$\vct a\cdot\vct b=ab\cos\theta$取沿一方向的投影, 得标量; 叉积$\vct a\times\vct b$大小为$ab\sin\theta$, 方向垂直两者平面, 以右手从第一矢量转向第二矢量决定. 后面力矩与磁力会用到这个方向, 可先以$\uv x\times\uv y=\uv z$核查. 矢量本身的大小记$|\vct a|$, 不等于各分量的代数和.

静止的两个点电荷$q_1,q_2$相距$r$时, 力沿连线. 用$\rr_1,\rr_2$表示位置, $\rr_{21}=\rr_2-\rr_1$, 则$q_1$对$q_2$的力为
\[
\vct F_{1\to2}=\kc q_1q_2\frac{\rr_{21}}{|\rr_{21}|^3},\qquad \kc=\frac1{4\pi\eps}\simeq8.99\times10^9\,\mathrm{N\,m^2/C^2}.
\]
两电荷同号时沿$\rr_{21}$排斥, 异号时方向反转. 库仑定律是实验规律, 不是由牛顿定律单独推导的结论. 点电荷模型适用于物体尺度远小于观测距离的情形; 对真实分布用许多小电荷的力相加. 真空中的静电力遵从矢量叠加, 因而不能先把各力的大小相加再选择方向.

电场把“某处会怎样影响电荷”与“放入什么电荷”分开. 在测试电荷很小而不显著扰动源分布的极限, 定义$\ee=\vct F/q_{\rm test}$, 单位为N/C, 也等于V/m. 负测试电荷受力沿$-\ee$. 对固定源电荷,
\[
\ee(\rr)=\kc\sum_i q_i\frac{\rr-\rr_i}{|\rr-\rr_i|^3},\qquad \vct F=q\ee.
\]
电场是空间每一点上的矢量函数, 不同于一根实际绳索. 场线的切线方向沿场, 疏密可表示强弱; 静电场线始于正电荷、终于负电荷或无穷远. 同一点若场非零只有一个方向, 所以场线不能交叉. 场线本身不是带电粒子的轨迹, 因为轨迹还依赖初速度与惯性.

\begin{center}
\begin{tikzpicture}[>=Stealth,scale=.85]
\filldraw[structurecolor] (-2,0) circle(.12) node[below=4pt] {$+q$};
\filldraw[main] (2,0) circle(.12) node[below=4pt] {$-q$};
\draw[->](-1.8,0)--(1.8,0);
\foreach \s in {-1,1}{\draw[->](-1.95,.18*\s) .. controls(-1,2*\s) and(1,2*\s)..(1.95,.18*\s);}
\node at(0,-2){电偶极子的部分场线(示意)};
\end{tikzpicture}
\end{center}

重力也可分为源和响应: 质量$M$建立$\vct g=-GM\uv r/r^2$, 质量$m$的受力为$m\vct g$. 电场与它结构相似, 但电荷有正负且同号排斥, 普通质量之间只有吸引. 温度$T(x,y,z)$则每点只给一个数, 是标量场. 等值线将同一温度连接, 色图用颜色标数值, 高度图用绘图高度标数值; 高度不是把原来平面物体真的抬高. 矢量箭头还表达方向与大小, 无箭头的纹理只表示局部轴向, 未必区分正反. 读图前须确认它编码的是哪种对象.

\originalsection{电偶极子与场中运动}
等量异号电荷相隔小矢量$\vct d$(由负指向正)形成电偶极子, 电偶极矩定义为$\vct p=q\vct d$, 单位C$\cdot$m. 在均匀外场中两力相抵, 合力为零, 但力矩为$\vct\tau=\vct p\times\ee$. 因为两力各自有力臂, 标量大小为$qEd\sin\theta=pE\sin\theta$. 磁针类比之前先看电偶极的转动. 从角$\theta$缓慢逆抗恢复力矩转动$\dd\theta$, 外功$\dd U=pE\sin\theta\dd\theta$, 积分$U=-pE\cos\theta+\text{常数}$. 因而可取$U=-\vct p\cdot\ee$, 对齐$\theta=0$为极小和稳定平衡, 反对齐$\theta=\pi$为极大、稍偏即转离. 非均匀场中两端场值不同, 中性偶极子也可能受净力; 这解释带电棒吸引中性纸屑时极化的作用.

在均匀场中, 忽略其他力且速度远低于光速, $m\ddot{\rr}=q\ee$给出匀加速运动. 若入射速度垂直场, 沿初速度方向匀速、沿场方向匀加速, 轨迹类似抛物线. 场为零处也不一定有静止粒子; 力为零只说明该瞬间加速度为零.
